\documentclass[10pt,a4paper]{amsart}
\usepackage{amsmath,amssymb,amsthm,graphicx}
\usepackage{verbatim}
\newenvironment{reference}{\comment}{\endcomment}
%\newenvironment{reference}{}{}
\newenvironment{slogan}{\comment}{\endcomment}
\newenvironment{history}{\comment}{\endcomment}

% For commutative diagrams we use Xy-pic
\usepackage[all]{xy}

% We use 2cell for 2-commutative diagrams.
\xyoption{2cell}
% Theorem environments.
%
\theoremstyle{plain}
\newtheorem{theorem}[subsection]{Theorem}
\newtheorem{proposition}[subsection]{Proposition}
\newtheorem{lemma}[subsection]{Lemma}

\theoremstyle{definition}
\newtheorem{definition}[subsection]{Definition}
\newtheorem{example}[subsection]{Example}
\newtheorem{exercise}[subsection]{Exercise}
\newtheorem{situation}[subsection]{Situation}

\theoremstyle{remark}
\newtheorem{remark}[subsection]{Remark}
\newtheorem{remarks}[subsection]{Remarks}

\numberwithin{equation}{subsection}

% Macros
%
\def\lim{\mathop{\mathrm{lim}}\nolimits}
\def\colim{\mathop{\mathrm{colim}}\nolimits}
\def\Spec{\mathop{\mathrm{Spec}}}
\def\Hom{\mathop{\mathrm{Hom}}\nolimits}
\def\Ext{\mathop{\mathrm{Ext}}\nolimits}
\def\SheafHom{\mathop{\mathcal{H}\!\mathit{om}}\nolimits}
\def\SheafExt{\mathop{\mathcal{E}\!\mathit{xt}}\nolimits}
\def\Sch{\mathit{Sch}}
\def\Mor{\mathop{\mathrm{Mor}}\nolimits}
\def\Ob{\mathop{\mathrm{Ob}}\nolimits}
\def\Sh{\mathop{\mathit{Sh}}\nolimits}
\def\NL{\mathop{N\!L}\nolimits}
\def\CH{\mathop{\mathrm{CH}}\nolimits}
\def\proetale{{pro\text{-}\acute{e}tale}}
\def\etale{{\acute{e}tale}}
\def\QCoh{\mathit{QCoh}}
\def\Ker{\mathop{\mathrm{Ker}}}
\def\Im{\mathop{\mathrm{Im}}}
\def\Coker{\mathop{\mathrm{Coker}}}
\def\Coim{\mathop{\mathrm{Coim}}}

% Boxtimes
%
\DeclareMathSymbol{\boxtimes}{\mathbin}{AMSa}{"02}

%
% Macros for moduli stacks/spaces
%
\def\QCohstack{\mathcal{QC}\!\mathit{oh}}
\def\Cohstack{\mathcal{C}\!\mathit{oh}}
\def\Spacesstack{\mathcal{S}\!\mathit{paces}}
\def\Quotfunctor{\mathrm{Quot}}
\def\Hilbfunctor{\mathrm{Hilb}}
\def\Curvesstack{\mathcal{C}\!\mathit{urves}}
\def\Polarizedstack{\mathcal{P}\!\mathit{olarized}}
\def\Complexesstack{\mathcal{C}\!\mathit{omplexes}}
% \Pic is the operator that assigns to X its picard group, usage \Pic(X)
% \Picardstack_{X/B} denotes the Picard stack of X over B
% \Picardfunctor_{X/B} denotes the Picard functor of X over B
\def\Pic{\mathop{\mathrm{Pic}}\nolimits}
\def\Picardstack{\mathcal{P}\!\mathit{ic}}
\def\Picardfunctor{\mathrm{Pic}}
\def\Deformationcategory{\mathcal{D}\!\mathit{ef}}
\title{Exactness of the associated graded forces exactness of the filtered complex}
\author{The Stacks Project Authors}
\date{stacks-project:homology.tex:local-snapshot-20261012; GNU FDL 1.2 or later}
\begin{document}
\maketitle
\noindent\textit{Editorial note.} All delivered mathematical text is the author's own, copied from the pinned Stacks Project chapter file homology.tex; no statement or proof step was rewritten, reordered or removed.\par\medskip
\noindent\textit{Editorial note.} The proof is the author's complete argument; 'The proof of (2) is dual.' and 'Of course (5) follows from (2).' are the author's own routine reductions to a part that is written out in full inside the counted span, and both sentences are preserved verbatim.\par\medskip
\noindent\textit{Editorial note.} Every labelled reference occurring inside the counted spans is carried in assembly.external\_references as a non-load-bearing navigation pointer (same-chapter ones at their definition site, cross-chapter ones at the invocation inside the counted span), so no load-bearing mathematics is replaced by a stub. No bibliography citation occurs inside the counted spans.\par\medskip
\noindent\textit{Editorial note.} The same-name check for this card scanned every staged Stacks chapter of this repository -- all sources/stacks-*/pinned-original-source.tex files, the card's own chapter excluded -- and the labels found are recorded in dedup.same\_name\_cross\_chapter on the card.\par\medskip
\noindent\textit{Editorial note.} The preamble is the chapter's own native preamble excerpt already used for this source: the comment-type reference, slogan and history environments (native preamble lines 8-13), the Xy-pic setup (native preamble lines 14-18) and the native theorem declarations with the mathematical macros (native preamble lines 154-217, the actual last line of the file).\par\medskip
\noindent\textit{Editorial note.} Same-name check: the label lemma-filtered-acyclic also occurs in the staged chapter stacks-derived (sources/stacks-derived/pinned-original-source.tex, line 4509). The two are not the same claim and do not double count: the derived-chapter lemma states that the filtered acyclic complexes form a saturated triangulated subcategory of the filtered homotopy category and identifies the associated multiplicative system, whereas this card counts the author's statement about a complex of filtered objects with finite filtrations whose associated graded is exact. The assumptions and the conclusions are different, so both may be counted by their respective chapters.\par\medskip
\noindent\textit{Allowed standard prerequisites (editorial).}\begin{itemize}
\item Filtered objects, finite filtrations and the associated graded. Same chapter, the definitions of a filtered object, its filtration steps F\textasciicircum{}p and the graded pieces gr\textasciicircum{}p. Supplies the notation of the statement and, in the induction, the decomposition of the associated graded of an object into the top filtration step plus the graded of the quotient.
\item Short exact sequences of filtered objects. Same chapter, the exact sequence of filtration steps and quotients induced by a short exact sequence of filtered objects. Supplies the commutative diagram with short exact rows used in the induction step.
\end{itemize}\medskip
\begin{lemma}
\label{lemma-filtered-acyclic}
Let $\mathcal{A}$ be an abelian category.
Let $A \to B \to C$ be a complex of filtered objects of $\mathcal{A}$.
Assume $A, B, C$ have finite filtrations and that
$\text{gr}(A) \to \text{gr}(B) \to \text{gr}(C)$ is exact.
Then
\begin{enumerate}
\item for each $p \in \mathbf{Z}$ the sequence
$\text{gr}^p(A) \to \text{gr}^p(B) \to \text{gr}^p(C)$ is exact,
\item for each $p \in \mathbf{Z}$ the sequence
$F^p(A) \to F^p(B) \to F^p(C)$ is exact,
\item for each $p \in \mathbf{Z}$ the sequence
$A/F^p(A) \to B/F^p(B) \to C/F^p(C)$ is exact,
\item the maps $A \to B$ and $B \to C$ are strict, and
\item $A \to B \to C$ is exact (as a sequence in $\mathcal{A}$).
\end{enumerate}
\end{lemma}

\begin{proof}
Part (1) is immediate from the definitions.
We will prove (3) by induction on the length of the filtrations.
If each of $A$, $B$, $C$ has only one
nonzero graded part, then (3) holds as $\text{gr}(A) = A$, etc.
Let $n$ be the largest integer such that at least one of
$F^nA, F^nB, F^nC$ is nonzero. Set $A' = A/F^nA$, $B' = B/F^nB$,
$C' = C/F^nC$ with induced filtrations. Note that
$\text{gr}(A) = F^nA \oplus \text{gr}(A')$
and similarly for $B$ and $C$. The induction hypothesis
applies to $A' \to B' \to C'$, which implies that
$A/F^p(A) \to B/F^p(B) \to C/F^p(C)$ is exact for $p \geq n$.
To conclude the same for $p = n + 1$, i.e., to prove that $A \to B \to C$
is exact we use the commutative diagram
$$
\xymatrix{
0 \ar[r] & F^nA \ar[r] \ar[d] & A \ar[r] \ar[d] & A' \ar[r] \ar[d] & 0 \\
0 \ar[r] & F^nB \ar[r] \ar[d] & B \ar[r] \ar[d] & B' \ar[r] \ar[d] & 0 \\
0 \ar[r] & F^nC \ar[r] & C \ar[r] & C' \ar[r] & 0
}
$$
whose rows are short exact sequences of objects of $\mathcal{A}$.
The proof of (2) is dual. Of course (5) follows from (2).

\medskip\noindent
To prove (4) denote $f : A \to B$ and $g : B \to C$ the given morphisms.
We know that $f(F^p(A)) = \Ker(F^p(B) \to F^p(C))$ by (2) and
$f(A) = \Ker(g)$ by (5). Hence
$f(F^p(A)) =  \Ker(F^p(B) \to F^p(C)) =
\Ker(g) \cap F^p(B) = f(A) \cap F^p(B)$ which proves that
$f$ is strict. The proof that $g$ is strict is dual to this.
\end{proof}

\end{document}
